\documentclass{article}

\usepackage[margin=1in]{geometry}
\usepackage{xcolor}
\usepackage{stepdiff}

\title{A Polished stepdiff Note}
\author{}
\date{}

\stepdiffsetup{
  display=notes,
  theme=soft,
  layout=lecture,
  highlight=background,
  reason-style=muted,
  color-mode=typed
}

\begin{document}

\maketitle

\section*{Sum of the first integers}

Pair the sum with a reversed copy. Typed highlights keep additions and final emphasis visually distinct while staying token-level.

\[
\begin{stepdiff}[frame=true]
\step{S_n = 1 + 2 + \cdots + (n - 1) + n}
\step[reason={reverse order}]{S_n = n + (n - 1) + \cdots + 2 + 1}
\step[reason={add the two lines}]{2S_n = (n + 1) + (n + 1) + \cdots + (n + 1)}
\step[reason={count n terms}]{2S_n = n(n + 1)}
\step[reason={divide by 2}, diff=all, tag=final]{S_n = \frac{n(n + 1)}{2}}
\end{stepdiff}
\]

\section*{A limit comparison}

\[
\begin{stepdiff}[display=focus, theme=focus, color-mode=teaching, reason-style=badge, frame=true, legend=true]
\step{a_n \le b_n}
\step[reason={take limits}]{\lim a_n \le \lim b_n}
\step[reason={identify limits}, diff=all, tag=final]{L \le M}
\end{stepdiff}
\]

\section*{A derivative example}

\[
\begin{stepdiff}[color-mode=typed]
\step{g(x) = x^2\sin x}
\step[reason={product rule}]{g'(x) = 2x\sin x + x^2\cos x}
\step[reason={evaluate at zero}, diff=all, tag=final]{g'(0) = 0}
\end{stepdiff}
\]

\end{document}
